\subsection{6.1 训练集群}\label{training-cluster}

该集群被设计为一个系统级优化问题。所需算力包络由规模定律估算、数据质量目标、token 预算和实际观测到的训练效率推导得出。随后，基础设施规划将该包络转化为数据中心电力与散热、GPU、CPU、NVLink 纵向扩展域、InfiniBand 横向扩展网络、存储、系统软件、诊断、调度、检查点与可观测性等环节中可用的 FLOPs。

规划目标并非仅理论峰值 FLOPs。长时间运行的训练任务会因 MFU 损失、检查点开销、重计算、硬件损耗、维护、验证负载和恢复缓冲区而损失可用容量。因此，我们针对可用训练容量进行优化：即所交付硬件中能够被调度、保持健康、以高 MFU 运行并在故障发生时快速恢复的部分。

MAI-Base-1 在包含 8K 张 GB200 GPU 的单一逻辑集群上完成训练，MAI-Thinking-1 的 RL 爬升则使用 4.6K 张 GB300 运行。通过将训练保持在同一代同构加速器上、处于已知良好的机架健康边界内、在稳定的调度器行为下运行，并贴近可预测的存储路径，这一选择降低了实验方差。更广泛的实验室环境还包括用于开发、验证、对比剖析和下一代系统上线的 H100、GB200 与 GB300 系统，但主训练运行优先保证局部性、拓扑稳定性和运维同构性。

有关我们集群环境的更多细节见附录 L。

\subsection{6.2 训练稳定性、确定性与有效产出指标}\label{training-stability-determinism-and-goodput-metrics}

训练稳定性通过集群使任务持续高效推进、从故障中快速恢复并在各重启路径上保持数值正确性的能力来衡量。我们同时追踪显式故障与静默效率损失。显式故障包括崩溃循环、节点故障、InfiniBand 与 NVLink 链路抖动和链路中断、内存溢出错误、Pod 终止、检查点停滞以及人工重新入队。静默效率损失包括 MFU 下降、重计算、冗长的启动路径、缓慢的 actor 调度、检查点引发的停滞、内存性能退化，以及不会立即导致任务崩溃但会降低有效吞吐量的网络状况。

确定性被视为一等基础设施属性。在前沿规模下，确定性训练不仅是模型代码层面的事，它还取决于完整执行基座（execution substrate）的正确性与稳定性。集群必须消除静默数据损坏（Hochschild et al., 2021; Dixit et al., 2021），保持通信拓扑稳定，防止不健康的设备与链路进入任务，并在检查点与重启边界之间保持浮点归约和累加的顺序。

有效产出（goodput）是首要的生产 KPI。我们将有效产出定义为理想训练时长与实际墙钟时长之比（Jiang et al., 2024; Grattafiori et al., 2024），并将两者差距分解为各类开销。这一框架使失败成本变得显式化。一次失败不仅耗费到重启为止的时间，还可能强制重算自上一个持久检查点以来的内容，引发启动与 actor 调度延迟，扰乱任务放置，阻塞检查点推进，或在恢复后导致 MFU 降低。同理，即使任务仍在运行，MFU 下降也会被当作生产事故处理，因为集群仍在消耗已分配的 GPU 小时数，而产出的训练进展低于预期。

\subsection{6.2.1 MFU 与有效产出}\label{mfu-and-goodput}

总体趋势是从可靠性向效率的转变。第一层改进来自减少中断，如崩溃循环、节点故障、链路抖动、重新入队、内存溢出（OOM）错误和检查点停滞。第二层、也更难的一层改进来自消除静默的有效产出损失，如 MFU 下降、重计算、冗长的启动路径、缓慢的进程调度、检查点引发的停滞，以及网络或内存性能退化。在前沿规模下，两层改进都很重要，因为每多一小时开销都会在上千张 GPU 上成倍放大。运营目标是将有效产出视为首要生产 KPI：每种故障模式、运行时变慢、检查点事件和 MFU 回退都有明确的负责人、检测信号、预防路径，以及对可用 FLOPs 的量化影响。

MAI-Base-1 预训练在 8K 张 GPU 上达到了 90.0\% 的有效产出，尽管其规模大于此前的预训练任务。总开销降至 51 小时。重计算——即回退到检查点后重新生成先前已计算步骤所花的时间——降至 6.5 小时，仅占总开销的 15\%。非推进时间降至 14 小时，占 27\%，表明系统在保持存活、避免重复返工以及在无需长时间人工干预的情况下恢复方面有了显著改善。不过，最终运行也清楚暴露了下一个瓶颈：MFU 下降开销成为剩余类别中最大的一项，达 18 小时，占总开销的 35\%。这主要来自检查点写入、网络退化、内存压力和硬件健康状态迁移。故障趋势虽有改善，但并未消失。

\subsection{6.3 推理效率与模型部署}\label{inference-efficiency-and-model-deployment}

随着推理模型日益重要，推理成本、延迟和服务可扩展性成为部署的关键约束。因此，我们在模型开发的整个过程中将推理效率视为一等目标，在技术栈各层寻找优化机会：模型架构、服务引擎，乃至部署硬件的选择。

为提升每瓦性能，我们在微软的 MAIA-200 硬件上实现 MAI-Thinking-1。与基于 GB200 的部署相比，在相同机架功率预算下，MAIA-200 的 token 生成吞吐量高出 40\% 以上。每瓦性能的提升使数据中心电力得到更高效的利用，并支持在更大规模上承载推理负载。

\subsection{6.4 可持续性与社区优先的 AI 基础设施倡议}\label{sustainability-and-communityfirst-ai-infrastructure-initiatives}

对我们而言，AI 的发展必须兼顾社会与环境可持续性，并惠及我们建设与运营基础设施所在的社区。微软承诺实现碳负排放、水正效益、零废弃，并保护生态系统。2025 年，公司在实现全球年用电量 100\% 由可再生能源匹配的目标上达成了一项里程碑。

MAI-Thinking-1 主要在微软自有的亚利桑那州凤凰城（Phoenix, AZ）基础设施上完成训练，并在得克萨斯州达拉斯（Dallas, TX）完成后续训练。凤凰城的设施按照 LEED 金级认证标准建设，以表彰其在环境可持续性和能源效率方面的卓越表现。例如，我们这里的备用发电机使用可再生柴油，与传统石油基柴油相比可降低净碳排放。

微软的 Community-First AI Infrastructure 倡议承诺最小化环境影响，并投资于其建设与运营所在社区。该方式以与凤凰城地区公用事业公司、水务部门和环保组织的紧密合作为基础，在降低整体系统需求的同时强化共享基础设施。

微软已承诺投入超过 $50M 用于扩大公共水资源资产，包括与古德伊尔市（City of Goodyear）合作建设市政蓄水设施，并为惠及社区的供水与污水管道延伸及更广泛的系统升级提供资金。这些投资与大自然保护协会（The Nature Conservancy）以及 Gila River 蓄水系统等合作伙伴共同开展的保护行动相配合，帮助补给地下水并改善流域层面的用水效率。

与此同时，微软通过设在埃斯特雷拉山社区学院（Estrella Mountain Community College）和格兰代尔社区学院（Glendale Community College）的数据中心学院（Datacenter Academy）投资长期人才培养，确保凤凰城基础设施的增长转化为社区持久的经济机遇。

\subsection{7 结论与未来方向}\label{conclusion-and-future-directions}

我们介绍了我们的爬山式优化机器——一种优化流水线每一个组成部分（从数据和基础设施到 RL 配方与评测）的模型开发方法。MAI-Thinking-1 是该机器产出的首个模型：一个 35B active / 1T total 参数的 MoE，其训练不依赖对第三方模型的蒸馏。在 STEM 推理与软件工程任务上，MAI-Thinking-1 位列其参数量级中最强模型之列。

MAI-Thinking-1 是起点而非终点。展望未来，我们计划将爬山式优化扩展到更多模态、更大规模和更精细的能力。AI 的进步不是任何单一模型的产物，而是可以被可靠改进的流水线的产物。我们将继续打磨我们的流水线，并期待分享未来的爬升进展。
